%% Allgemeines -----------------------------------

% Fortlaufende Nummerierung - Entweder 1, 2, 3, 4, ... oder mit einem semantischen Schema
\def \decisionID {1}

% Datum, an dem der Record anefertigt wurde
\def \decisionDate {22.11.2026}

% Wer hat an dem Record mitgewirkt?
\def \decisionParticipants {Niklas Käsehagen, OneTutor (Customer Proxy)}

% Der Status in dem sich die Anforderung in diesem Record aktuell befindet.
% Möglichkeiten: Vorgeschlagen, Akzeptiert, Abgelehnt, Abgelöst
\def \decisionStatus {Abgelöst}

% ID des Records, der diesen Record ablöst, falls zutreffend
\def \decisionSuperseded {2}

% ID des Records, den dieser Record ablöst, falls zutreffend
\def \decisionSuperseding {}

%% Anforderung -----------------------------------

% Titel der Anforderung
\def \requirementTitle {Folien dürfen nicht ”gefressen” werden}

% Kontext der Anforderung
\def \requirementContext {Wenn Folien in das Repository hochgeladen werden, entscheidet die Software Similar Slides, ob die Folie
einzeln, neben ähnlichen, sich bereits im Repository befindlichen Folien, oder nicht gespeichert werden soll. Dabei darf es nicht passieren, dass eine Folie, die einzigartig ist, fälschlicherweise verworfen wird, obwohl sie sich noch nicht im Repository befindet.}

% Konsequenz der Anforderung
\def \requirementConsequences {Das Ähnlichkeitsmaß muss genau erkennen können, ob eine Folie im Repository enthalten ist oder nicht und eine Folie nur nicht speichern, wenn sicher ist, dass sie sich genau so bereits im Repository befindet.}